\section{Coefficient rings, minuscule classes, and flavor}
\label{sec:preliminaries}

This section supplies the generation and coefficient maps needed for the
two ideal comparisons described in Section~\ref{sec:introduction}.  Minuscule
generation makes $B_F$ a quotient of $R$; Kirwan surjectivity does the same
for $H_F^*(X)$.  Flavor base change and local constancy of regular
stability will let us apply these presentations to the auxiliary data
used later.

We use complex coefficients throughout.  Vertices with $v_i=0$ can be
discarded.  If all $v_i$ vanish, both sides of Theorem~\ref{thm:main} equal
$H_F^*(\mathrm{pt})$, so we henceforth allow ourselves to omit that case.
Write
\[
 P_F=\C[x_{ia}:i\in Q_0,\ 1\le a\le v_i]\otimes\C[\mathfrak f],
 \qquad \mathcal W=\prod_i S_{v_i},\qquad R=P_F^{\mathcal W}.
\]
The $x_{ia}$ are the Chern roots of the gauge bundles.  Each linear
form in $P_F$ has cohomological degree two.  An element
$\lambda=(\lambda_{ia})\in X_*(T)$ has topological charge
$[\lambda]=(\sum_a\lambda_{ia})_i$.
We write $A_F^{(d)}$ for the subspace of charge $d$ and define
\begin{equation}\label{eq:full-fixed-quotient}
 B_F=A_F/\langle A_F^{(d)}:d\ne0\rangle.
\end{equation}
Thus $B_F$ is the coordinate ring of the fixed subscheme for
$T_{\mathrm{top}}$.  A subscript $0$ means that the flavor torus is trivial.

\begin{lemma}\label{lem:trivial-summands}
For the purposes of Theorem~\ref{thm:main}, one may assume that the flavor
action assigns a character to each arrow summand and acts diagonally on
each framing multiplicity space.  Adding or removing a $G$-trivial
representation with any flavor action preserves both coefficient-algebra
quotients and all stability conditions.
\end{lemma}

\begin{proof}
Decompose $N$ into $G$-isotypic summands.  A commuting torus acts on each
multiplicity space and is diagonalizable there.  The nontrivial irreducible
gauge representations that occur are the bifundamentals between distinct
vertices, the dual standard representations from framing, and the
traceless endomorphisms at vertices of rank at least two.  All remaining
summands are gauge-trivial.  A weighted traceless endomorphism summand can
be replaced by the full endomorphism representation by adding a scalar
line of the same flavor weight.  This gives the asserted description.

If $E$ is gauge-trivial, the moment map for $N\oplus E$ is the moment map
for $N$, and semi-invariant generation shows that the stable and
semistable loci are the corresponding loci for $N$, times $T^*E$.
Their stabilizers are unchanged.  The quotient gains a factor $T^*E$,
which retracts $F$-equivariantly to zero.  In the space of triples the
condition on the $E(\OO)$ coordinate imposes no lattice constraint, so
this coordinate gives a common vector-space factor in every finite
truncation.  The normalized equivariant Thom isomorphism removes this
factor, compatibly with convolution, both gradings, and the coefficient
map.  This proves the claim on both sides.
\end{proof}

We call a coweight \emph{positive minuscule} if all its entries, across
all vertices, are $0$ or $1$, and \emph{negative minuscule} if all entries
are $0$ or $-1$.  Its orbit $\Gr_G^\lambda$ in
$\Gr_G=G(\KK)/G(\OO)$ is a product of ordinary Grassmannians and is
closed and smooth.

Write $\mathcal W_\lambda=\operatorname{Stab}_{\mathcal W}(\lambda)$,
$u_+=\max(u,0)$, and let $\wt(N)$ denote the multiset of weights of $N$.
For $p\in P_F^{\mathcal W_\lambda}$, put $p_{w\lambda}=w(p)$; stabilizer
invariance makes this independent of the representative $w$.

\begin{proposition}[Abelianization and minuscule generation]
\label{prop:minuscule}
There is an injective coefficient-preserving algebra map from $A_F$ into
\[
 \bigl(\operatorname{Frac}(P_F)[X_*(T)]\bigr)^{\mathcal W},\qquad
 e^\lambda e^\gamma=e^{\lambda+\gamma}.
\]
Besides $R$, the algebra is generated by the dressed classes associated
with all nonzero positive and negative minuscule coweights.  For
$p\in P_F^{\mathcal W_\lambda}$, the class has expression
\begin{equation}\label{eq:minuscule}
 \sum_{\lambda'\in\mathcal W\lambda}
 e^{\lambda'}p_{\lambda'}
 \frac{\prod_{\beta\in\wt(N)}
           \beta^{(-\langle\beta,\lambda'\rangle)_+}}
      {\eu_T(T_{t^{\lambda'}}\Gr_G^\lambda)}.
\end{equation}
Each distinct translate occurs once, and the dressing is translated
with the coweight.  Denote this class by $M_{\lambda,p}$.
Weights $\beta$ include their flavor components, but
pairings use only the gauge components; factors with exponent zero are
omitted.  The class has charge $[\lambda]$.  On charge zero, the BFN
grading is twice the ordinary polynomial degree of the rational
coefficients in this description.
\end{proposition}

\begin{proof}
We use the classical BFN localization and flavor construction
\cite[Proposition~3.22, Lemmas~5.9--5.13, Proposition~5.15,
and Corollary~5.21]{BFN18}, and its minuscule formula
\cite[Proposition~6.6]{BFN18}\cite[Appendix~A, Proposition~A.2]{BFN19},
together with quiver minuscule generation
\cite[Proposition~3.1 and Remarks~3.3--3.4]{Weekes19}.
The form of the formula can be checked directly: the inclusion of the
space of triples into the bundle of lattices $gN(\OO)$ contributes the
Euler class of
\[
 t^{\lambda'}N(\OO)\big/
       \bigl(N(\OO)\cap t^{\lambda'}N(\OO)\bigr),
\]
which is the numerator of \eqref{eq:minuscule}.  Localization on the
closed minuscule orbit gives the tangent denominator.  There is no
rotation, so translating a modification by $t^\eta$ changes the Laurent
monomial but does not shift the coefficients.  The localization theorem
and the orbit filtration give injectivity before taking Weyl invariants.
Equivalently, the image of the nonabelian algebra lies in the invariant
part of the displayed ring.

For completeness, the minuscule coweights used in the generation theorem
are subsets of slots at several vertices at once.  Decomposing positive
or negative integral coweights by successive level sets gives exactly
these two families; no assertion that single-slot monopoles suffice is
needed.  Lemma~\ref{lem:trivial-summands} puts arbitrary commuting flavors
in weighted quiver form.  Flavor characters change the equivariant Euler
factors but not the gauge-weight hyperplane arrangement used in the
generation proof; the same chamber argument therefore retains those
parameters, as discussed in \cite[Section~1.6.1]{Weekes19}.
The embedding used here is the classical localization map, not a
specialization of a representation in which loop rotation is inverted.
Finally, put $c_\lambda=\sum_{\beta\in\wt(N)}
(-\langle\beta,\lambda\rangle)_+$ and
$a_\lambda=\dim\Gr_G^\lambda$.  For homogeneous $p$, the BFN homological
normalization gives
\[
 \deg M_{\lambda,p}=2c_\lambda-2a_\lambda+\deg p
\]
by \cite[Lemmas~2.2 and 2.5--2.6]{BFN18}.  These are exactly the degrees
of the numerator, tangent denominator, and dressing in
Equation~\eqref{eq:minuscule}.  The bare matter-free torus class
$e^\lambda$ has degree zero.  This proves the stated grading comparison.
General dressings are decomposed into homogeneous summands.
The additional shifts from a common finite section cutoff will be kept
explicitly in Section~\ref{sec:correspondences}.
No nonnegative grading on $A_F$ itself is assumed.
\end{proof}

Every nontrivial minuscule generator in Proposition~\ref{prop:minuscule} has nonzero
charge.  Consequently the coefficient map $R\to B_F$ is surjective,
and $B_F$ is nonnegatively graded.  More precisely, the ideal in
\eqref{eq:full-fixed-quotient} is generated by those minuscule classes:
after killing them the algebra is generated by $R$ and is supported in
charge zero.  In particular, a coefficient relation in $B_F$ can be
expressed as a finite sum of neutral words of positive length in the
nonzero minuscule generators, with coefficients in $R$.

\begin{lemma}[Flavor base change]\label{lem:flavor-base-change}
Let $F\to F'$ be a homomorphism of algebraic tori, and suppose the action
on $N$ is obtained by restriction of an $F'$-action commuting with $G$.
Then the natural maps are isomorphisms
\begin{align*}
 A_{F'}\otimes_{H_{F'}^*(\mathrm{pt})}H_F^*(\mathrm{pt})&\ \cong A_F,\\
 B_{F'}\otimes_{H_{F'}^*(\mathrm{pt})}H_F^*(\mathrm{pt})&\ \cong B_F.
\end{align*}
They respect the coefficient maps, both gradings, and convolution.
\end{lemma}

\begin{proof}
Filter the affine Grassmannian by bounded unions of orbits and truncate
all section lattices between sufficiently high positive and negative
powers of $t$.  Over an orbit, the truncated space of triples is a
vector bundle of constant relative rank.  The orbit is an affine-space
bundle over a partial flag variety; its equivariant homology, with the
relative shifts in the definition of $A_F$, is even.  Its $F'$-equivariant
homology is free over $H_{F'}^*(\mathrm{pt})$, and restricting the flavor
action gives the displayed change of coefficients on this stratum.
The exact sequences for adjoining strata have no odd terms, so these
statements extend through each finite filtration.  The free successive
quotients also show that tensoring the sequences introduces no Tor term.
Passing to the stabilized truncations and their exhaustive union proves
the first assertion.  Pullback, refined intersection, and proper
pushforward defining convolution commute with this change of
equivariance, so the assertion is an algebra isomorphism.

Flavor parameters have charge zero.  Under the first isomorphism the
ideal generated by all nonzero charge spaces is the extension of the
same ideal before base change.  Taking the quotient therefore gives
the second assertion; no flatness of $B_{F'}$ is required for this step.
\end{proof}

For the moment-map comparisons below, we use the topological affine
Kempf--Ness theorem: the quotient of a character-shifted real moment
level by a compact form of $G$ is homeomorphic to the corresponding
affine-character GIT quotient
\cite[Theorem~4.2(v)]{Hoskins14}.  The same statement restricts to a closed
invariant subvariety.  Indeed, its GIT quotient embeds as a closed
subspace of the ambient quotient, and its orbit closures and their
closed representatives remain in that subvariety.  This applies in
particular to every central complex moment fiber in $T^*N$.

\begin{lemma}[Local constancy of regular stability]\label{lem:local-stability}
Let $\theta$ be a rational character parameter for which every semistable
point of $\mu^{-1}(0)$ is stable and has trivial stabilizer.  There is a
neighborhood $U$ of $\theta$ in $X^*(G)\otimes_{\Z}\mathbb R$ such that
the stable and semistable loci are unchanged for every rational
$\theta'\in U$.  One may choose such a $\theta'$ outside all proper
hyperplanes $\theta'\cdot\mathbf b=0$ with $0<\mathbf b\leq\mathbf v$,
and then clear its denominators by a positive integer.  Empty loci are included.
\end{lemma}

\begin{proof}
We identify a positive multiple of a rational character with the same
stability condition.  Fix a maximal torus of $G$ and the finite set of its
weights on $N\oplus N^*$.  For a point $x$ and $g\in G$, let $C(gx)$ be the
cone generated by the weights occurring in $gx$.  The affine
Hilbert--Mumford criterion says that $x$ is semistable precisely when the
character belongs to every such cone, and is stable with finite stabilizer
precisely when it belongs to all their full interiors in the real weight
space; see \cite{King94} and
\cite[Propositions~2.5 and 2.7]{Hoskins14}.
Although $g$ varies, only finitely many cones
occur.  A sufficiently small neighborhood of $\theta$ remains inside each
cone whose interior contains $\theta$, and outside each cone that does not
contain $\theta$.  A $\theta$-stable point therefore remains stable, and a
$\theta$-unstable point remains unstable.  Regularity accounts for all
points of the closed subvariety $\mu^{-1}(0)$.  This proves the assertion
about the entire loci, including when they are empty.  The geometric
quotient and its principal gauge bundle are unchanged, as is freeness.
Choose a rational point in this neighborhood avoiding the finitely many
proper hyperplanes specified in the statement, and then rescale by a
positive integer.  If $G$ is trivial, its character space is a point and
the list of bounded nonzero dimension vectors is empty, so the assertion
is immediate.
\end{proof}

The Kirwan map $R\to H_F^*(X)$ is surjective.  For nonempty $X$,
Lemma~\ref{lem:local-stability} permits a rational perturbation $\theta$
of $\nu$ that preserves the entire stable locus $D$ and avoids
$\theta\cdot\mathbf b=0$ for $0<\mathbf b\leq\mathbf v$.  Its extension
$(\theta,-\theta\cdot\mathbf v)$ is nondegenerate for the framed
dimension vector $(\mathbf v,1)$: a proper subvector has final
coordinate zero, or its nonzero complement does.  Clearing denominators
preserves stability.  The quotient, its principal gauge bundle, and its
tautological bundles are unchanged, so ordinary Kirwan surjectivity
applies with the same coefficient map
\cite[Definition~3.1 and Theorem~1.2]{MN18}; see also
\cite[Section~2.1]{DK25}.  The torus-equivariant conclusion is
\cite[Corollary~1.5]{MN18}; we recall its argument for an arbitrary
commuting flavor torus.  Choose a homogeneous basis of $H^*(X)$ represented by polynomials
in the ordinary tautological Chern classes.  These classes lift to
$F$-equivariant Chern classes.  Leray--Hirsch for
$X\to EF\times_F X\to BF$ therefore makes $H_F^*(X)$ free over
$H_F^*(\mathrm{pt})$ on these lifts and proves the asserted equivariant
surjectivity.  In particular $H^*(X)$ is finite-dimensional and even.
For empty $X$ the assertion means the unique surjection to the zero ring.
